% ═══════════════════════════════════════════════════════════════
% MUSTERLÖSUNG - DS-01: Mi nueva casa
% Klasse 9, Unidad 2
% ═══════════════════════════════════════════════════════════════

\documentclass[10pt, a4paper]{article}

\usepackage[utf8]{inputenc}
\usepackage[T1]{fontenc}
\usepackage[ngerman]{babel}

\usepackage{palatino}
\usepackage{helvet}
\newcommand{\sanstext}[1]{{\fontfamily{phv}\selectfont #1}}

\usepackage{microtype}
\usepackage[margin=1.4cm, top=1.6cm, bottom=1.3cm]{geometry}
\usepackage{enumitem}
\usepackage{tabularx}
\usepackage{booktabs}
\usepackage{array}
\usepackage{multicol}
\usepackage{tikz}
\usetikzlibrary{calc}
\usepackage{amssymb}

\usepackage{fancyhdr}
\usepackage{lastpage}
\usepackage{needspace}

\usepackage{xcolor}
\usepackage{tcolorbox}
\tcbuselibrary{skins}

% ═══════════════════════════════════════════════════════════════
% FARBEN
% ═══════════════════════════════════════════════════════════════

\definecolor{kopfzeile}{HTML}{1A1A1A}
\definecolor{akzent}{HTML}{C62828}
\definecolor{hellgrau}{HTML}{F2F2F2}
\definecolor{mittelgrau}{HTML}{777777}
\definecolor{dunkelgrau}{HTML}{444444}
\definecolor{loesung}{HTML}{C62828}  % Rot für Lösungen

% ═══════════════════════════════════════════════════════════════
% VARIABLEN
% ═══════════════════════════════════════════════════════════════

\newcommand{\thema}{Mi nueva casa}
\newcommand{\klasse}{9}
\newcommand{\maxpunkte}{30}

% ═══════════════════════════════════════════════════════════════
% KOPF- UND FUSSZEILE
% ═══════════════════════════════════════════════════════════════

\pagestyle{fancy}
\fancyhf{}
\renewcommand{\headrulewidth}{0pt}
\renewcommand{\footrulewidth}{0pt}
\cfoot{\sanstext{\footnotesize\textcolor{mittelgrau}{\thepage\,/\,\pageref{LastPage}}}}

% ═══════════════════════════════════════════════════════════════
% BOXEN
% ═══════════════════════════════════════════════════════════════

\newtcolorbox{grammatikbox}{
    colback=hellgrau,
    colframe=hellgrau,
    boxrule=0pt,
    arc=0pt,
    left=8pt, right=8pt, top=8pt, bottom=6pt,
    borderline north={1.5pt}{0pt}{dunkelgrau}
}

\newtcolorbox{vokabelbox}{
    colback=white,
    colframe=mittelgrau,
    boxrule=0.5pt,
    arc=0pt,
    left=6pt, right=6pt, top=6pt, bottom=6pt
}

\newcommand{\regel}[1]{%
    \tikz[baseline=(X.base)]\node[fill=akzent, text=white,
    font=\sanstext\scriptsize\bfseries, inner sep=2pt, rounded corners=1pt](X){#1};%
}

% ═══════════════════════════════════════════════════════════════
% BEFEHLE
% ═══════════════════════════════════════════════════════════════

\newcommand{\es}[1]{\textbf{#1}}
\newcommand{\de}[1]{{\small\textit{\textcolor{mittelgrau}{#1}}}}
\newcommand{\luecke}[1]{\rule{#1}{0.4pt}}
\newcommand{\punktebox}[1]{\hfill\sanstext{\small[\_\_/#1]}}

% Lösung in Rot
\newcommand{\lsg}[1]{\textcolor{loesung}{\textbf{#1}}}

\newcommand{\aufgabe}[2]{%
    \needspace{4\baselineskip}%
    \vspace{10pt}%
    \noindent\sanstext{\textbf{#1}\quad#2}%
    \vspace{5pt}%
    \nopagebreak[4]%
}

\newlist{teil}{enumerate}{1}
\setlist[teil]{label=\alph*), leftmargin=1.8em, itemsep=3pt, parsep=0pt, topsep=3pt}

\setlength{\parindent}{0pt}
\setlength{\parskip}{4pt}
\setlength{\columnsep}{20pt}

% ═══════════════════════════════════════════════════════════════
% DOKUMENT
% ═══════════════════════════════════════════════════════════════

\begin{document}

% ═══════════════════════════════════════════════════════════════
% HEADER
% ═══════════════════════════════════════════════════════════════

\noindent
\begin{tikzpicture}[remember picture, overlay]
    \fill[kopfzeile] (current page.north west) rectangle +(\paperwidth, -1cm);
    \node[anchor=west, text=white, font=\sanstext\large\bfseries]
        at ([xshift=1.4cm, yshift=-0.5cm]current page.north west) {\thema{} -- \textcolor{loesung}{MUSTERLÖSUNG}};
    \node[anchor=east, text=white, font=\sanstext\small]
        at ([xshift=-1.4cm, yshift=-0.5cm]current page.north east)
        {Spanisch\,·\,Klasse \klasse\,·\,ML};
\end{tikzpicture}

\vspace{0.5cm}

\noindent
\sanstext{\small\textbf{Name:}} \luecke{5cm} \hfill
\sanstext{\small\textbf{Datum:}} \luecke{2.5cm}

\vspace{6pt}

% ═══════════════════════════════════════════════════════════════
% KASTEN G1: BEWEGUNGSVERBEN (MIT LÖSUNGEN)
% ═══════════════════════════════════════════════════════════════

\begin{grammatikbox}
\sanstext{\footnotesize\textbf{Kasten G1: Los verbos de movimiento}}

\vspace{4pt}

\begin{multicols}{2}

{\small Ergänze die Präpositionen:}

\vspace{3pt}

{\footnotesize
\begin{tabular}{@{}cll@{}}
$\downarrow$ & \es{bajar} \de{herunterbringen} & \lsg{de / del} \\[0.5em]
$\uparrow$ & \es{subir} \de{hochbringen} & \lsg{a / al} \\[0.5em]
$\oplus$ & \es{poner} \de{stellen, legen} & \lsg{en} \\
\end{tabular}}

\columnbreak

{\footnotesize
\begin{tabular}{@{}cll@{}}
$\ominus$ & \es{sacar} \de{herausnehmen} & \lsg{de / del} \\[0.5em]
$\leftarrow$ & \es{traer} \de{bringen (her)} & \lsg{de / del} \\[0.5em]
$\rightarrow$ & \es{llevar} \de{bringen (hin)} & \lsg{a / al} \\
\end{tabular}}

\end{multicols}

\end{grammatikbox}

\vspace{4pt}

% ═══════════════════════════════════════════════════════════════
% KASTEN G2: OBJEKTPRONOMEN (MIT LÖSUNGEN)
% ═══════════════════════════════════════════════════════════════

\begin{grammatikbox}
\sanstext{\footnotesize\textbf{Kasten G2: Los pronombres de objeto directo}}

\vspace{4pt}

{\small Ergänze die Pronomen:}

\vspace{3pt}

{\footnotesize
\begin{tabular}{@{}llll@{}}
\textbf{Singular} & & \textbf{Plural} & \\
\lsg{lo} \de{ihn} & $\leftarrow$ maskulin $\rightarrow$ & \lsg{los} \de{sie (m)} \\[0.4em]
\lsg{la} \de{sie (f)} & $\leftarrow$ feminin $\rightarrow$ & \lsg{las} \de{sie (f)} \\
\end{tabular}}

\end{grammatikbox}

\vspace{6pt}

% ═══════════════════════════════════════════════════════════════
% AUFGABEN MIT LÖSUNGEN
% ═══════════════════════════════════════════════════════════════

\aufgabe{Aufgabe 1}{Ordne die Bilder den Vokabeln zu.} \punktebox{6}

\vspace{2pt}

\begin{multicols}{2}
{\footnotesize
\begin{tabular}{@{}cl@{}}
1. $\square$ Bett & \lsg{b)} \\
2. $\square$ Lampe & \lsg{d)} \\
3. $\square$ Kühlschrank & \lsg{f)} \\
4. $\square$ Sofa & \lsg{c)} \\
5. $\square$ Schrank & \lsg{a)} \\
6. $\square$ Fernseher & \lsg{e)} \\
\end{tabular}

\columnbreak

\begin{tabular}{@{}cl@{}}
7. $\square$ Badezimmer & \lsg{j)} \\
8. $\square$ Küche & \lsg{i)} \\
9. $\square$ Wohnzimmer & \lsg{g)} \\
10. $\square$ Schlafzimmer & \lsg{h)} \\
11. $\square$ Flur & \lsg{l)} \\
12. $\square$ Balkon & \lsg{k)} \\
\end{tabular}
}
\end{multicols}

\vspace{8pt}

% ═══════════════════════════════════════════════════════════════

\aufgabe{Aufgabe 2}{Ergänze die Verben und Präpositionen.} \punktebox{8}

\vspace{4pt}

\begin{teil}
    \item Los chicos de la mudanza \lsg{bajan} los muebles \lsg{del} camión.

    \item Después \lsg{suben} las cajas \lsg{al} nuevo piso.

    \item Clara \lsg{pone} el póster \lsg{en} la pared.

    \item Mamá \lsg{saca} la ropa \lsg{de} la caja.
\end{teil}

\vspace{8pt}

% ═══════════════════════════════════════════════════════════════

\aufgabe{Aufgabe 3}{Ersetze das Nomen durch ein Pronomen.} \punktebox{6}

\vspace{4pt}

\begin{teil}
    \item \es{Ponemos la lámpara en el salón.} $\rightarrow$ \lsg{La ponemos en el salón.}
    \item \es{Subimos los muebles al piso.} $\rightarrow$ \lsg{Los subimos al piso.}
    \item \es{Llevamos el sofá al dormitorio.} $\rightarrow$ \lsg{Lo llevamos al dormitorio.}
    \item \es{Sacamos las plantas del camión.} $\rightarrow$ \lsg{Las sacamos del camión.}
    \item \es{Bajamos la cama del piso.} $\rightarrow$ \lsg{La bajamos del piso.}
    \item \es{Traemos el televisor del salón.} $\rightarrow$ \lsg{Lo traemos del salón.}
\end{teil}

\vspace{8pt}

% ═══════════════════════════════════════════════════════════════

\aufgabe{Aufgabe 4}{Beschreibt gemeinsam einen Umzug.} \punktebox{10}

\vspace{4pt}

{\footnotesize\textbf{Mögliche Dialoge:}}

\vspace{2pt}

{\small
\textbf{Partner A:} \es{¿Dónde ponemos \lsg{el sofá}?}\\
\textbf{Partner B:} \lsg{Lo ponemos en el salón.}

\vspace{3pt}

\textbf{Partner A:} \es{¿Y dónde está \lsg{la nevera}?}\\
\textbf{Partner B:} \lsg{Está en la cocina.}

\vspace{3pt}

\textbf{Partner B:} \es{¿Dónde ponemos \lsg{la cama}?}\\
\textbf{Partner A:} \lsg{La ponemos en el dormitorio.}

\vspace{3pt}

\textbf{Partner B:} \es{¿Qué hay en \lsg{el baño}?}\\
\textbf{Partner A:} \lsg{Hay una ducha y un espejo.}
}

% ═══════════════════════════════════════════════════════════════
% FOOTER
% ═══════════════════════════════════════════════════════════════

\vfill

\noindent\rule{\textwidth}{0.5pt}

\vspace{4pt}

\hfill\sanstext{\textbf{Punkte:}} \luecke{1.2cm} / \maxpunkte

\end{document}
